% -------------------------------------------------------
\section{Analytical Framework and Hypotheses}
% -------------------------------------------------------

The analysis is organized around two main parts: how the text is represented and how changes are compared. Together, these determine which types of changes can be observed and how reliably they are detected. Using different ways to represent and compare the same set of texts helps assess reliability by checking whether different methods agree. When results are supported by multiple methods, they are more trustworthy than those from a single method.\\

\textbf{Axis 1 -- Text Representation} \newline
This axis concerns how abstracts are encoded. Lexical representations (term frequency, TF-IDF) capture surface vocabulary change and provide interpretability, but cannot detect semantic shift when wording remains stable. Topic models (LDA, NMF) infer latent thematic structure from co-occurrence patterns, enabling detection of broader shifts in research focus, but introduce modelling assumptions and require manual interpretation. \\ Embedding-based representations (Word2Vec, SentenceBERT, SPECTER2) capture contextual and semantic similarity, allowing detection of conceptual and discourse-level change even when vocabulary is stable, but are less interpretable and require triangulation. Together, these representations separate lexical change (what is said), thematic change (what is discussed), and semantic change (how ideas are framed).\\
This axis improves performance by allowing comparison between sparse lexical signals and richer semantic representations, reducing reliance on surface vocabulary and increasing sensitivity to conceptual change.\newline

\textbf{Axis 2 -- Comparison Method} \newline
This axis concerns how temporal change is measured.
Direct tracking measures change using counts, topic shares, or centroid movement, which is easy to understand but cannot show how features work together. Classification methods identify key differences across time periods but can confuse topic changes and how ideas are presented. Geometric methods (PCA, UMAP) reveal patterns in the embedding space; PCA shows overall differences, while UMAP preserves local groupings, but both primarily describe the data rather than explain it.\\
Centroid-based drift provides a way to measure change over time that works across different representations, allowing direct comparison between TF-IDF and embedding spaces.\\
This axis makes results more reliable by combining easy-to-understand tracking methods with geometric and predictive approaches, letting us assess time-based changes through both direct measures and hidden patterns.\\

\textbf{Sensitivity, Robustness, and the Role of Convergence}
Methods differ systematically in sensitivity: lexical approaches capture vocabulary change, embeddings capture semantic shift, topic models capture thematic structure, and classification captures separability but conflates multiple signals. This heterogeneity is used as a design feature rather than a limitation: combining methods reduces dependence on any single representational bias and improves robustness to method-specific noise. As a result, convergence across independent approaches provides stronger evidence of genuine temporal change, while divergence is treated as informative about different layers of discourse evolution.


\noindent Four hypotheses structure the evaluation:

\begin{description}[leftmargin=0pt]
\item[H1:] Topic models recover the lexical shift from theory to application while revealing transitional hybrid themes.
\item[H2:] Embedding-based methods detect conceptual shifts earlier than TF-IDF due to changes in semantic context preceding vocabulary adoption.
\item[H3:] Centroid drift captures fine-grained temporal transitions, while classification reflects broader separability rather than precise inflection points.
\item[H4:] Major temporal patterns will be broadly consistent across representations and distance metrics, indicating robustness to methodological choice, while also exhibiting systematic differences in sensitivity.
\end{description}
